\documentclass[a4paper,12pt]{article}




%\abovedisplayshortskip=-12pt plus 3pt
%\setlength{\abovedisplayskip}{-5pt}
%\setlength{\belowdisplayskip}{-5pt}
%\linespread{1.05}
%\voffset=-14cm
%\hoffset=-0.5cm

\usepackage{bm, amsmath, amsfonts, amssymb, url}
\usepackage{titlesec, titletoc}
\usepackage{graphicx,subfigure}
\usepackage{appendix, indentfirst, makeidx, array}
\usepackage{booktabs}
\usepackage{fancyhdr}
\pagestyle{fancy}
\usepackage{fullpage} 

\addtolength{\headsep}{+1cm}
\addtolength{\oddsidemargin}{+0.5cm}
\addtolength{\evensidemargin}{-0.5cm}
%\usepackage[cam,a4,center]{crop}
%\usepackage{geometry}


%\linespread{1.3}

%\fancypagestyle{plain}{
%\fancyhf{}
%\renewcommand{\headrulewidth}{0pt}
%%\fancyfoot[c]{\thepage}
%}


%\addtocontents{toc}{\protect\thispagestyle{empty}}

%%%-------page style ------------

%\setlength{\baselineskip}{20pt plus0.3pt minus0.2pt}
%\setlength{\marginparsep}{0pt}
%\setlength{\marginparwidth}{1.5cm}

%\setlength{\oddsidemargin}{3cm}
%\setlength{\evensidemargin}{3cm}
%%%-------------------------------


\renewcommand{\labelenumi}{(\alph{enumi})}
\renewcommand{\arraystretch}{1.2}
\titleformat{\chapter}[display]{\Large\bf}{\LARGE{Chapter~\thechapter}}{2mm}{}
\titleformat{\section}{\large\bf}{\thesection}{1.5mm}{}
\titleformat{\subsection}{\bfseries\itshape}{\thesubsection}{1.5mm}{}
%self-define symbol
\newcommand{\proposition}[1]{\bigskip\noindent\fbox{\begin{minipage}{0.95\textwidth}~\begin{minipage}{0.95\textwidth}#1\end{minipage}\\\smallskip\end{minipage}}\bigskip\vline\vline\vline}

\newcommand{\excercise}{\addtocounter{section}{+1}
\subsection*{\underline{Ex~\thesection}}}
\newcommand{\stag}[1]{\tag{\arabic{chapter}.\arabic{equation}#1}}

\newcommand{\qdo}[0]{\quad}
\newcommand{\qdt}[0]{\quad\quad}
\newcommand{\qdthr}[0]{\quad\quad\quad}
\newcommand{\qdf}[0]{\quad\quad\quad\quad}
\newcommand{\qdfiv}[0]{\quad\quad\quad\quad\quad}

\newcommand{\norm}[1]{\|#1\|}
%\newcommand{\sdelta}{{\scriptstyle\Delta}}
\newcommand{\sdelta}{\Delta}
\newcommand{\bra}[1]{\langle#1|}
\newcommand{\ket}[1]{|#1\rangle }
\newcommand{\bket}[1]{\Big |#1\Big\rangle }
\newcommand{\braket}[2]{\langle#1|#2\rangle }
\newcommand{\proj}[2]{|#1\rangle\langle#2|}
\newcommand{\expect}[2]{\langle #1\rangle_{#2}}
\newcommand{\Kappa}[0]{\bm{\mathcal{K}}}
\newcommand{\parity}[0]{\bm{\mathcal{P}}}
\newcommand{\parityl}[0]{\bm{p}}
\newcommand{\charge}[0]{\bm{\mathcal{C}}}


\makeindex
%-----------------------------------%

\begin{document}

\abovedisplayskip=-6pt plus 3pt minus 0pt
\belowdisplayskip=8pt plus 3pt minus 0pt

\abovedisplayshortskip=-6pt plus 3pt minus 0pt
\belowdisplayshortskip=8pt plus 3pt minus 0pt

\title{Problem Set 7 (Due: Nov. 9, 2011)}
\maketitle

Your homework should be submitted to TA (Chih-Han Lin) at the  beginning of the course. Generally, You have a week to finish the problem set since it was revealed. If you cannot submit it on time, you can still bring your homework to TA's office (L325, Institute of Atomic and Molecular Sciences, in NTU campus) and put it into TA's mailbox which is on the top of the shoe cabinet in L325. Always remember to remind TA to check your homework by e-mail (\url{clin@ltl.iams.sinica.edu.tw}) if you don't submit it in class.

%\abovedisplayskip=0pt %plus 3pt minus 3pt
%\belowdisplayskip=2pt %plus 3pt minus 3pt
%
%\abovedisplayshortskip=0pt %plus 3pt minus 3pt
%\belowdisplayshortskip=2pt %plus 3pt minus 3pt

\excercise
For the system of a linear chain with $N$ mass point of $m$ equal spacing in the lengh $l$. Each mass point is attached by a first spring of force constant $\alpha$ to the nearest neighbor mass points, and by the second spring of force constant $\beta$ to the equilibrium positions $x_r^0=rl/N=ra$. Then the Lagrangian of the system can be expressed as

\begin{align}
\mathcal{L}=\frac{m}{2}\sum_r\dot{x_r}^2-\frac{\alpha}{2}\sum\left(
x_{r+1}-x_r-a\right)^2-\frac{\beta}{2}\sum_r(x_r-x_r^0)^2.
\end{align}

\begin{figure}
\centering
\includegraphics[width=0.7\textwidth]{oscillator.pdf} 
\caption{figure of ex.1}
\end{figure}

\begin{enumerate}
\item Define the displace of $r$-th mass point from its equilibrium position by $q_r$, then

\begin{align}
\mathcal{L}=\frac{m}{2}\sum_r\dot{q_r}^2-\frac{\alpha}{2}\sum\left(
q_{r+1}-q_r\right)^2-\frac{\beta}{2}\sum_r(q_r)^2.
\end{align}

\item Define $\varphi(x_r^0, t)$ as a new dynamical variable such that $q_r(t)=\sqrt{a/m}\varphi (x_r^0, t)$ then show that the action of the system $A$ can be written as

\begin{align}
A=\int dt dx \frac{1}{2}\left\{\left(\frac{\partial\varphi}{\partial t}\right)^2-
\frac{T}{\mu} \left(\frac{\partial\varphi}{\partial x}\right)^2+\frac{T}{\mu}\frac{1}{\lambda_c}\varphi^2\right\}.
\end{align}

in the limit $\alpha, N\longrightarrow \infty$ and $\beta, m \longrightarrow 0$ such that $\alpha a\longrightarrow T$, $m/a\longrightarrow \mu$ and

\begin{align*}
\frac{\beta a}{m}\longrightarrow \frac{T}{\mu}\frac{1}{\lambda_c^2}.
\end{align*}

where $A=\int dt dx \mathcal{L}(\varphi, \dot{\varphi}, \varphi')$ with

\begin{align*}
\mathcal{L}=\frac{1}{2}(\dot{\varphi}-\frac{T}{\mu}\varphi'^2-\frac{T}{\mu}\frac{1}{\lambda_c^2}\varphi^2).
\end{align*}
\end{enumerate}

\excercise

Verify that

\begin{align}
\sum_r(\bm{q}_r-\bm{q}_{r+1})^2&=\frac{\hbar}{m}\sum_s 4\sin^2\left(\frac{2\pi s}{N}\right)\bm{Q}_s\bm{Q}_s^\dag,
\end{align}

if 

\begin{align*}
\bm{q}_r=\sqrt{\frac{\hbar}{mN}}\sum_s\exp\left(\frac{2\pi i}{N}sr\right)\bm{Q}_s.
\end{align*}

\excercise

Verify that

\begin{align*}
[\bm{a}_{s} , \bm{a}_{s'} ]&=[\bm{a}_{s}^\dag , \bm{a}^{\dag}_{s'}]=0,\\
[\bm{a}_{s}, \bm{a}^\dag_{s'}]&=\delta_{ss'}\mathbf{I}.
\end{align*}


\end{document}
